\scanpage{384}\mspage{123}

d) Il s'ensuit que les hom (*) sont des \uline{iso} de $\mathrm{Hot}$, donc tous les $C_\infty(\underline t)$ sont isomorphes dans $\underline{\mathrm{Hot}}$. L'astuce du $t*t'$ (dans le cas où on a une loi $t*t'$ associative et admettant l'objet final comme \inserted{objet} unité) donne $\simeq$ des isom \uline{canoniques}

\[
C_\infty(t)\xrightarrow{\sim} C_\infty(t')\qquad\text{dans }\mathrm{Hot}
\]

et c'est là un système transitif d'iso. Donc on a dans $\mathrm{Hot}$ un objet \deleted{$C_\infty$, défini} canonique

\[
C_\infty\in\mathrm{Ob}\,\mathrm{Hot},
\]

dont la relation à $H\in\mathrm{Ob}\,\mathrm{Ind\,Hot}$ est déduite de c), c'est une flèche canonique

\[
H\longrightarrow C_\infty\qquad\text{dans }\mathrm{Ind\,Hot}.
\]

% slanted
\margin{\uline{NB} $H$ n'est pratiquement jamais ess. constant, donc la flèche \deleted{\unsure{préci|préc.}} ci-contre n'est pas un iso dans $\mathrm{Ind\,Hot}$.}

Cette flèche induit un iso. pour les objets correspondants de $\widehat{\mathrm{Hot}_0}$\,; \unsure{Quand} on regarde $H$ \inserted{non} comme \deleted{fonct.} objet de $\widehat{\mathrm{Hot}_0}$, mais prenant le foncteur \inserted{$H_0$} induit sur $\mathrm{Hot}_0$ par le foncteur ind-représentable $H$ sur $\mathrm{Hot}$. Ce foncteur induit est ``représentable'' par l'objet \inserted{$C_\infty$} de la catégorie plus grande

\[
H_0\simeq C_\infty\qquad\text{dans }\widehat{\mathrm{Hot}_0},
\]

(si le foncteur \inserted{can.} $\mathrm{Hot}\longrightarrow\widehat{\mathrm{Hot}_0}$ est pl. fid, alors cette formule décrit $C_\infty$ \deleted{\unsure{comme}} objet de $\mathrm{Hot}$, à isom unique près, en termes de $H$, \deleted{donc} \inserted{aussi} en termes de $H\in\mathrm{Ob}\,\mathrm{Ind\,Cat\,hot}$.

% archive sheet number: 383
